\section{Foundations: application}
\begin{frame}[plain]{Foundations: empirical application}
\Large Islamic political representation and women's education\par
\vspace{1em}\normalsize
The sharp RD illustration from Foundations, after the theoretical discussion.\\[1em]
\textbf{Stata module 1} and worksheet A.
\src{Original study: \textcite{meyersson2014}. Official companion code: CIT\_2020\_CUP\_polecon.do.}
\end{frame}

\begin{frame}{The Turkish municipal election design}
\begin{items}
\item \textbf{Units:} 2,629 municipalities in the authors' teaching dataset.
\item \textbf{Score:} Islamic party margin of victory in the 1994 mayoral election, measured in percentage points.
\item \textbf{Treatment:} an Islamic party wins the mayoralty, $T_i=\ind\{X_i\geq0\}$.
\item \textbf{Outcome:} percentage of women aged 15--20 who completed high school, measured in 2000.
\item \textbf{Target:} effect of an Islamic mayor for municipalities at a tied election.
\end{items}
\src{\textcite{meyersson2014}; Foundations, Section 2.1.}
\end{frame}

\begin{frame}{The raw outcome-score relationship}
\fig[.73]{turkey_raw.pdf}
\src{Stata module 1, authors' dataset. A comparison of all winners and losers mixes the cutoff effect with the outcome-score relationship.}
\end{frame}

\begin{frame}{Evenly spaced bins in the application}
\fig[.73]{turkey_bins_es.pdf}
\src{Twenty bins on each side, separate linear descriptive fits. Stata module 1; Foundations, Section 3.}
\end{frame}

\begin{frame}{Quantile spaced bins in the application}
\fig[.73]{turkey_bins_qs.pdf}
\src{Twenty bins on each side, separate linear descriptive fits. Different bin widths reflect the uneven score density.}
\end{frame}

\begin{frame}[fragile]{Local estimation in Stata}
\begin{lstlisting}
* An interaction allows different slopes on the two sides.
generate double TX = T*X
regress Y T X TX if abs(X)<=20, vce(robust)

* Data driven local estimation and robust inference.
rdrobust Y X, p(1) kernel(triangular) bwselect(mserd) all
rdrobust Y X, p(1) kernel(triangular) bwselect(cerrd) all
\end{lstlisting}
\begin{items}
\item $Y$ is already a percentage. A coefficient of 3 means 3 percentage points.
\item The robust RD interval accounts for estimated smoothing bias.
\end{items}
\F{Sections 4.2--4.3; module 1}
\end{frame}

\begin{frame}{The local fit around a tied election}
\fig[.73]{turkey_local.pdf}
\src{Module 1. Local linear, triangular weights, MSE selected bandwidth. The plot uses the estimation sample and bandwidth.}
\end{frame}

\begin{frame}{Point estimates and robust intervals}
\begin{center}
\small\renewcommand{\arraystretch}{1.3}
\begin{tabular}{@{}lrrr@{}}\toprule
 & Estimate & Bias-corrected & Robust 95\% CI\\\midrule
MSE bandwidth & 3.02 & 2.98 & $[-0.31, 6.28]$\\
CER bandwidth & 2.43 & 2.41 & $[-1.16, 5.98]$\\
\bottomrule\end{tabular}
\end{center}

\begin{items}
\item Units: percentage points of female high-school completion.
\item The MSE point estimate is about $+3.02$ pp. Its robust interval includes zero.
\item The data remain compatible with both a small negative effect and a substantively positive effect.
\end{items}
\src{Saved Stata output, module 1. Current package results, rather than copied book values.}
\end{frame}

\begin{frame}{Predetermined covariates}
\begin{center}
\small\renewcommand{\arraystretch}{1.3}
\begin{tabular}{@{}lrrr@{}}\toprule
 & Estimate & Bias-corrected & Robust 95\% CI\\\midrule
Log population, 1994 & -0.08 & -0.08 & $[-0.81, 0.65]$\\
Number of parties, 1994 & -0.26 & -0.29 & $[-1.58, 1.00]$\\
Islamic mayor, 1989 & 0.08 & 0.09 & $[-0.07, 0.26]$\\
\bottomrule\end{tabular}
\end{center}

\begin{items}
\item Each row retains the covariate's own units.
\item These checks use pre-treatment population, party count, and the previous mayor.
\item Variables measured in 2000 can respond to the 1994 treatment. Their timing requires separate justification.
\end{items}
\src{Module 1, coverage error selected bandwidths. A robust interval containing zero does not demonstrate exact balance.}
\end{frame}

\begin{frame}{Density evidence at the election threshold}
\fig[.73]{turkey_density.pdf}
\src{Module 1, rddensity. Interpret the formal test together with election rules and opportunities for precise sorting.}
\end{frame}

\begin{frame}{Bandwidth and exclusion sensitivity}
\begin{center}
\small\renewcommand{\arraystretch}{1.3}
\begin{tabular}{@{}lrrr@{}}\toprule
 & Estimate & Bias-corrected & Robust 95\% CI\\\midrule
Half bandwidth & 1.80 & 1.26 & $[-3.28, 5.81]$\\
Baseline & 3.02 & 2.98 & $[-0.31, 6.28]$\\
1.5 times bandwidth & 2.97 & 3.19 & $[0.32, 6.07]$\\
Local quadratic & 2.77 & 2.66 & $[-1.28, 6.60]$\\
Donut & 3.75 & 3.85 & $[0.41, 7.29]$\\
\bottomrule\end{tabular}
\end{center}

\begin{items}
\item The first three rows vary both $h$ and $b$ proportionally.
\item The donut drops $|X|\leq0.5$ and selects new bandwidths.
\item A significant donut estimate does not override uncertainty in the full analysis.
\end{items}
\src{Module 1. Outcome units are percentage points. Local quadratic uses its own selected bandwidth.}
\end{frame}

\begin{frame}{Sharp RD exercise and discussion}
\begin{enumerate}\setlength{\itemsep}{.6em}
\item Explain why the full-sample mean difference is not the cutoff estimand.
\item Compare the two binning rules with the local fit.
\item Identify the point estimate and the center of the robust interval.
\item Evaluate the timing of each proposed balance covariate.
\item Write a conclusion that reports magnitude, uncertainty, and policy scope.
\end{enumerate}
\src{Worksheet A, module 1. Allow 35--45 minutes for running and interpreting the analysis.}
\end{frame}
